%% ========================================================================
%% Preamble: Package dan Konfigurasi
%% ========================================================================

%% ------------------------------------------------------------------------
%% Encoding dan Font
%% ------------------------------------------------------------------------
\usepackage[utf8]{inputenc}      % Input encoding UTF-8
\usepackage[T1]{fontenc}         % Font encoding
%% ------------------------------------------------------------------------
%% Font Utama - Times New Roman 11pt
%% ------------------------------------------------------------------------
\usepackage{times}                  % Times New Roman font
\usepackage{amsfonts}               % Additional math fonts

%% Monospace font configuration for better readability
\usepackage[scaled=1.05]{inconsolata}  % Better monospace font, slightly enlarged
\usepackage{microtype}           % Typographical refinements
\usepackage[htt]{hyphenat}       % Allow hyphenation of dictionary words set in \texttt

%% ------------------------------------------------------------------------
%% Bahasa
%% ------------------------------------------------------------------------
\usepackage[english]{babel}  % Dukungan bahasa
\babelprovide[import]{indonesian}

%% Language flags - MUST be set by book-id.tex or book-en.tex before loading preamble
%% \newif\ifIndonesian
%% \newif\ifEnglish
%% Then set \Indonesiantrue or \Englishtrue
%%
%% The flags control all language-specific elements below

%% Konfigurasi istilah bahasa - conditional based on language
\providecommand{\abstractname}{Abstract}
\ifIndonesian
    \addto\captionsindonesian{%
        \renewcommand{\figurename}{Gambar}%
        \renewcommand{\tablename}{Tabel}%
        \renewcommand{\contentsname}{Daftar Isi}%
        \renewcommand{\listfigurename}{Daftar Gambar}%
        \renewcommand{\listtablename}{Daftar Tabel}%
        \renewcommand{\bibname}{Daftar Pustaka}%
        \renewcommand{\indexname}{Indeks}%
        \renewcommand{\appendixname}{Lampiran}%
        \renewcommand{\abstractname}{Abstrak}%
        \renewcommand{\partname}{Bagian}%
        \renewcommand{\chaptername}{Bab}%
    }
\else
    % English uses babel defaults
    \addto\captionsenglish{%
        \renewcommand{\bibname}{References}%
    }
\fi

%% ------------------------------------------------------------------------
%% Ukuran Kertas dan Margin
%% ------------------------------------------------------------------------
\usepackage{geometry}
\geometry{
    top=2cm,
    bottom=2cm,
    left=2cm,
    right=2cm,
    headheight=15pt
}

%% ------------------------------------------------------------------------
%% Grafis dan Gambar
%% ------------------------------------------------------------------------
\usepackage{graphicx}            % Untuk memasukkan gambar
\usepackage{float}               % Kontrol posisi float
\usepackage{subfig}              % Subfigure support
\usepackage{wrapfig}             % Bungkus teks di sekitar gambar
\raggedbottom                    % Hindari ruang vertikal besar saat relayout halaman

%% Float tuning for tighter page layout with tables/figures
\setcounter{topnumber}{3}
\setcounter{bottomnumber}{2}
\setcounter{totalnumber}{4}
\renewcommand{\topfraction}{0.9}
\renewcommand{\bottomfraction}{0.8}
\renewcommand{\textfraction}{0.07}
\renewcommand{\floatpagefraction}{0.8}

%% Konfigurasi caption - language-aware
\ifIndonesian
    \captionsetup{
        figurename=Gambar,
        tablename=Tabel
    }
\else
    \captionsetup{
        figurename=Figure,
        tablename=Table
    }
\fi

\usepackage{tikz}                % Untuk diagram
\usetikzlibrary{shapes,arrows,arrows.meta,positioning,calc,fit,backgrounds}
\usepackage{pgfplots}            % Plot dan grafik
\pgfplotsset{compat=1.18}

%% ------------------------------------------------------------------------
%% Gaya diagram TikZ (dipakai lintas bab; lihat docs/authoring-guide.md)
%% ------------------------------------------------------------------------
%% \diagbox: kotak netral abu-abu, dipakai untuk elemen sekunder pada diagram.
%% \diagbox accent: kotak beraksen warna bab (mis. blue!55!black untuk Bab 1),
%%   dipakai untuk elemen utama/penekanan pada diagram bab itu.
%% \diagarrow: panah alur, tebal, konsisten di semua diagram.
%% Seluruh tikzpicture pembahasan bab sebaiknya dibungkus
%% \resizebox{\linewidth}{!}{...} supaya ukurannya otomatis pas dengan lebar
%% teks halaman (11.5cm), tanpa perlu menghitung dimensi node secara manual.
\tikzset{
  diagbox/.style={
    draw=black!55, thick, rounded corners=2pt,
    fill=black!4,
    minimum width=2.1cm, minimum height=0.9cm,
    align=center, font=\small,
  },
  diagbox accent/.style={
    diagbox, draw=blue!55!black, fill=blue!8!white,
  },
  diagarrow/.style={
    -{Stealth[length=2.2mm]}, thick, black!65,
  },
}

%% Path untuk gambar
\graphicspath{{images/}{figures/}}

%% ------------------------------------------------------------------------
%% Tabel
%% ------------------------------------------------------------------------
\usepackage{booktabs}            % Tabel profesional
\usepackage{longtable}           % Tabel panjang (multi-halaman)
\usepackage{tabularx}            % Tabel dengan lebar kolom fleksibel
\usepackage{multirow}            % Sel tabel multi-baris
\usepackage{array}               % Enhanced array and tabular
% tabularray hooks into xcolor to load ninecolors, which uses rgb:Hsb color
% expressions. xcolor 2.x (Ubuntu's TeX Live) doesn't register the Hsb model
% by default, so ninecolors fails. We don't use ninecolors palettes, so stub
% it out before tabularray loads. xcolor is already loaded transitively by
% tikz/pgfplots above, so tabularray's hook fires immediately at registration.
\makeatletter
\@namedef{ver@ninecolors.sty}{9999/01/01 stubbed}
\makeatother
\usepackage{tabularray}           % Tabels dengan tblr environment
\usepackage{xstring}              % Ekstraksi substring (dipakai \branchref)

%% ------------------------------------------------------------------------
%% Warna
%% ------------------------------------------------------------------------
\usepackage{xcolor}              % Dukungan warna
\definecolor{codegreen}{gray}{0.25}
\definecolor{codegray}{gray}{0.45}
\definecolor{codepurple}{gray}{0.15}
\definecolor{backcolour}{gray}{0.94}
\definecolor{linkcolor}{gray}{0.0}
\definecolor{commandcolor}{gray}{0.0}

%% ------------------------------------------------------------------------
%% Listing Kode
%% ------------------------------------------------------------------------
\usepackage{listings}            % Syntax highlighting untuk kode
\usepackage{needspace}           % Reserve vertical space to avoid bad breaks

%% Konfigurasi nama caption untuk listing - language-aware
\ifIndonesian
    \renewcommand{\lstlistingname}{Kode}
    \renewcommand{\lstlistlistingname}{Daftar Kode}
\else
    \renewcommand{\lstlistingname}{Listing}
    \renewcommand{\lstlistlistingname}{List of Listings}
\fi

%% Definisi bahasa untuk listing: bash (perintah CLI)
\lstdefinelanguage{bash}{
    keywords={if, then, else, elif, fi, for, while, do, done, case, esac,
              function, return, exit, break, continue, in, select},
    sensitive=true,
    morecomment=[l]{\#},
    morestring=[b]",
    morestring=[b]'
}

%% Definisi bahasa untuk listing: Dart (bahasa utama buku ini).
%% listings tidak menyertakan definisi Dart bawaan, jadi didefinisikan manual.
\lstdefinelanguage{Dart}{
    morekeywords={abstract, as, assert, async, await, break, case, catch,
        class, const, continue, covariant, default, deferred, do, dynamic,
        else, enum, export, extends, extension, external, factory, false,
        final, finally, for, Function, get, hide, if, implements, import, in,
        interface, is, late, library, mixin, new, null, on, operator, part,
        required, rethrow, return, set, show, static, super, switch, sync,
        this, throw, true, try, typedef, var, void, while, with, yield},
    morekeywords=[2]{int, double, num, String, bool, List, Map, Set, Future,
        Stream, Widget, StatelessWidget, StatefulWidget, State, BuildContext,
        Column, Row, Container, Text, Scaffold, Navigator, Provider, override},
    sensitive=true,
    morecomment=[l]{//},
    morecomment=[s]{/*}{*/},
    morestring=[b]",
    morestring=[b]'
}

%% Basic style definition (used by both lstlisting and tcolorbox)
\lstdefinestyle{codestyle}{
    backgroundcolor=\color{backcolour},
    commentstyle=\color{codegreen}\itshape,
    keywordstyle=\color{black}\bfseries,
    keywordstyle=[2]\color{codepurple}\bfseries,
    numberstyle=\tiny\color{codegray},
    stringstyle=\color{codepurple},
    basicstyle=\ttfamily\scriptsize,
    breakatwhitespace=false,
    breaklines=true,
    captionpos=t,
    keepspaces=true,
    numbers=none,            % No line numbers
    showspaces=false,
    showstringspaces=false,
    showtabs=false,
    tabsize=2,
    frame=single,            % Simple frame around code
    rulecolor=\color{black},
    xleftmargin=0.6em,
    framexleftmargin=0.5em,
    framesep=2pt
}

%% Default language is PHP (the book's primary language, built into
%% listings); shell snippets pass [language=bash] explicitly.
\lstset{style=codestyle, language=PHP}

%% ------------------------------------------------------------------------
%% Prevent Page Breaks in Listings - Simple approach
%% ------------------------------------------------------------------------
%% Add penalties to discourage page breaks inside listings
\lstset{
    aboveskip=2pt,
    belowskip=0pt,
    abovecaptionskip=2pt,
    belowcaptionskip=0pt
}

%% Reserve 10 lines of vertical space before each lstlisting. Listings of
%% fewer than 10 lines start on a new page if needed and stay together;
%% longer listings break naturally past the reserved space.
\AddToHook{env/lstlisting/before}{\Needspace{10\baselineskip}}

%% For very long code that MUST be allowed to break, use this environment
\lstnewenvironment{longlisting}[1][]{%
    \lstset{#1}%
}{}

%% ------------------------------------------------------------------------
%% Hyperlinks dan Referensi
%% ------------------------------------------------------------------------
\usepackage[
    colorlinks=true,
    linkcolor=black,
    urlcolor=blue,
    citecolor=black,
    bookmarks=true,
    bookmarksnumbered=true,
    pdfborder={0 0 0}
]{hyperref}

%% ------------------------------------------------------------------------
%% Landscape pages (for wide tables that don't fit portrait text width)
%% ------------------------------------------------------------------------
\usepackage{pdflscape}

%% ------------------------------------------------------------------------
%% Header dan Footer
%% ------------------------------------------------------------------------
\usepackage{scrlayer-scrpage}    % KOMA-Script header/footer
\clearpairofpagestyles
\setkomafont{pagefoot}{\fontsize{10}{12}\selectfont}
\setkomafont{pagehead}{\fontsize{10}{12}\selectfont}
\ohead{\rightmark}               % Section in outer header
\ifoot{\leftmark}                % Chapter in inner footer
\ofoot[\pagemark]{\pagemark}    % Page number in outer footer (right on odd, left on even)
\pagestyle{scrheadings}

%% ------------------------------------------------------------------------
%% Bibliography (Daftar Pustaka)
%% ------------------------------------------------------------------------
\usepackage[
    backend=biber,
    style=ieee,
    sorting=none
]{biblatex}
\addbibresource{references.bib}
% biblatex has no indonesian.lbx; map it to english to suppress the
% "Language 'indonesian' not supported" warning that causes run.xml to
% oscillate and makes latexmk rerun pdflatex indefinitely.
\DeclareLanguageMapping{indonesian}{english}

%% ------------------------------------------------------------------------
%% Matematika
%% ------------------------------------------------------------------------
\usepackage{amsmath}             % Matematika lanjutan
\usepackage{amssymb}             % Simbol matematika
\usepackage{amsthm}              % Theorem environments

%% cleveref MUST be loaded after amsmath and hyperref
\usepackage[capitalise,noabbrev]{cleveref}  % Referensi pintar

%% Konfigurasi cleveref - language-aware
\ifIndonesian
    \crefname{figure}{gambar}{gambar}
    \Crefname{figure}{Gambar}{Gambar}
    \crefname{table}{tabel}{tabel}
    \Crefname{table}{Tabel}{Tabel}
    \crefname{chapter}{bab}{bab}
    \Crefname{chapter}{Bab}{Bab}
    \crefname{section}{bagian}{bagian}
    \Crefname{section}{Bagian}{Bagian}
    \crefname{subsection}{subbagian}{subbagian}
    \Crefname{subsection}{Subbagian}{Subbagian}
    \crefname{equation}{persamaan}{persamaan}
    \Crefname{equation}{Persamaan}{Persamaan}
    \crefname{listing}{kode}{kode}
    \Crefname{listing}{Kode}{Kode}
    \crefname{appendix}{lampiran}{lampiran}
    \Crefname{appendix}{Lampiran}{Lampiran}
\else
    % English defaults from cleveref are fine
    % Override only if needed
\fi

%% ------------------------------------------------------------------------
%% Box dan Environment Khusus
%% ------------------------------------------------------------------------
\usepackage[most]{tcolorbox}     % Colored boxes
\selectcolormodel{gray}         % Pastikan output aman untuk cetak hitam-putih
\usepackage{fontawesome5}        % Icons for box titles

%% Define box titles based on language
\ifIndonesian
    \def\noteboxTitle{\faInfoCircle\quad Catatan}
    \def\warningboxTitle{\faExclamationTriangle\quad Peringatan}
    \def\tipboxTitle{\faLightbulb\quad Tips}
    \def\exampleboxPrefix{\faBookOpen\quad Contoh: }
    \def\exerciseboxPrefix{Latihan }
    \def\challengeboxTitle{\faChevronCircleRight\quad Tantangan}
    \def\istilahpentingTitle{\faKey\quad Istilah Penting}
\else
    \def\noteboxTitle{\faInfoCircle\quad Note}
    \def\warningboxTitle{\faExclamationTriangle\quad Warning}
    \def\tipboxTitle{\faLightbulb\quad Tip}
    \def\exampleboxPrefix{\faBookOpen\quad Example: }
    \def\exerciseboxPrefix{Exercise }
    \def\challengeboxTitle{\faChevronCircleRight\quad Challenge}
    \def\istilahpentingTitle{\faKey\quad Key Terms}
\fi

% Box untuk catatan / note
\newtcolorbox{notebox}{
    colback=blue!5!white,
    colframe=blue!75!black,
    title=\noteboxTitle,
    fonttitle=\bfseries,
    enhanced, breakable,
    top=2pt, bottom=2pt, left=3pt, right=3pt,
    toptitle=1pt, bottomtitle=1pt,
    before skip=4pt, after skip=4pt
}

% Box untuk peringatan / warning
\newtcolorbox{warningbox}{
    colback=red!5!white,
    colframe=red!75!black,
    title=\warningboxTitle,
    fonttitle=\bfseries,
    enhanced, breakable,
    top=2pt, bottom=2pt, left=3pt, right=3pt,
    toptitle=1pt, bottomtitle=1pt,
    before skip=4pt, after skip=4pt
}

% Box untuk tip
\newtcolorbox{tipbox}{
    colback=green!5!white,
    colframe=green!75!black,
    title=\tipboxTitle,
    fonttitle=\bfseries,
    enhanced, breakable,
    top=2pt, bottom=2pt, left=3pt, right=3pt,
    toptitle=1pt, bottomtitle=1pt,
    before skip=4pt, after skip=4pt
}

% Box untuk contoh / example
\newtcolorbox{examplebox}[1][]{
    colback=yellow!10!white,
    colframe=orange!75!black,
    title=\exampleboxPrefix#1,
    fonttitle=\bfseries,
    enhanced, breakable,
    top=2pt, bottom=2pt, left=3pt, right=3pt,
    toptitle=1pt, bottomtitle=1pt,
    before skip=4pt, after skip=4pt
}

% Box untuk latihan / exercise
\newtcolorbox{exercisebox}[1][]{
    colback=purple!5!white,
    colframe=purple!75!black,
    title=\exerciseboxPrefix#1,
    fonttitle=\bfseries,
    enhanced, breakable,
    top=2pt, bottom=2pt, left=3pt, right=3pt,
    toptitle=1pt, bottomtitle=1pt,
    before skip=4pt, after skip=4pt
}

% Box untuk tantangan / challenge
\newtcolorbox{challengebox}{
    colback=teal!5!white,
    colframe=teal!75!black,
    title=\challengeboxTitle,
    fonttitle=\bfseries,
    enhanced, breakable,
    top=2pt, bottom=2pt, left=3pt, right=3pt,
    toptitle=1pt, bottomtitle=1pt,
    before skip=4pt, after skip=4pt
}

% Box untuk istilah penting / key terms (mengikuti preseden SQA: kotak ke-7 di
% luar katalog os-book, dipakai karena kosakata Dart/Flutter/Supabase padat)
\newtcolorbox{istilahpenting}{
    colback=green!5!white,
    colframe=green!75!black,
    title=\istilahpentingTitle,
    fonttitle=\bfseries,
    enhanced, breakable,
    top=2pt, bottom=2pt, left=3pt, right=3pt,
    toptitle=1pt, bottomtitle=1pt,
    before skip=4pt, after skip=4pt
}

%% Unbreakable code listing using tcolorbox (prevents page breaks)
%% This is an alternative to standard lstlisting for code that must stay together
\newtcblisting{codebox}[2][]{
    listing only,
    listing options={
        language=Dart,
        style=codestyle,
        float=,         % Disable float for tcolorbox
        #1
    },
    title={#2},
    colback=backcolour,
    colframe=black!30,
    fonttitle=\bfseries\small,
    breakable=false,    % This prevents page breaks inside the box
    enhanced,
    before upper={\parindent0pt}
}

%% For very long code that NEEDS to break across pages
\newtcblisting{longcodebox}[2][]{
    listing only,
    listing options={
        language=Dart,
        style=codestyle,
        float=,
        #1
    },
    title={#2},
    colback=backcolour,
    colframe=black!30,
    fonttitle=\bfseries\small,
    breakable=true,     % Allows page breaks for very long code
    enhanced,
    before upper={\parindent0pt}
}

%% ------------------------------------------------------------------------
%% Command Khusus Laravel/PHP (perintah CLI, class, paket, berkas, branch GitHub)
%% ------------------------------------------------------------------------
\newcommand{\cmd}[1]{\texttt{\textcolor{commandcolor}{#1}}}   % Perintah CLI (composer/php)
\newcommand{\phpclass}[1]{\texttt{#1}}                        % Nama class/model/controller PHP
\newcommand{\pkg}[1]{\textsf{#1}}                             % Nama paket dari Composer/npm
\newcommand{\file}[1]{\texttt{#1}}                            % Nama berkas konfigurasi/proyek
\newcommand{\artisan}[1]{\cmd{#1}}                            % Perintah Artisan CLI (php artisan ...)
\newcommand{\route}[1]{\texttt{#1}}                           % Path/nama route Laravel
\newcommand{\lstfile}[1]{\normalfont\ttfamily\scriptsize\textcolor{codegray}{\faIcon{file-alt}\ #1}} % Label path berkas di atas lstlisting

%% Repositori kode contoh, dua bentuk:
%%   - \repobase: repo multi-branch (dhanifudin/simple-pos), satu branch per bab,
%%     dipakai untuk narasi "baca riwayat commit berurutan" (Bab 1, Bab 12).
%%   - \branchref{}: repo template per-bab (dhanifudin/simple-pos-chNN), satu repo
%%     berdiri sendiri per bab dengan tombol "Use this template" GitHub, dipakai
%%     di akhir setiap Praktik supaya pembaca bisa mulai langsung dari bab mana pun.
\newcommand{\repobase}{dhanifudin/simple-pos}
\newcommand{\repoowner}{dhanifudin}

%% \branchref{chapter-08} -> callout baku "Kode lengkap: <repo template>, dari bab
%% chapter-08" dengan tautan langsung ke repo template GitHub-nya. Nomor bab (dua
%% digit terakhir dari argumen, mis. "08" dari "chapter-08") diekstrak lewat xstring
%% sehingga nama repo template (simple-pos-ch08) tidak perlu tabel manual per bab
%% dan otomatis tetap benar kalau suatu saat ada bab ke-13.
%% "branch" sengaja tidak diterjemahkan menjadi "cabang" di edisi Indonesia;
%% ini istilah teknis git/GitHub, diperlakukan sama seperti "widget"/"state".
%% Dipakai di akhir setiap Praktik; lihat docs/authoring-guide.md.
\newcommand{\branchref}[1]{%
    \StrRight{#1}{2}[\branchrefnum]%
    \def\branchreftemplate{simple-pos-ch\branchrefnum}%
    \begin{tcolorbox}[
        enhanced, breakable,
        colback=black!3, colframe=black!45,
        boxrule=0.5pt, arc=1.5pt,
        left=6pt, right=6pt, top=3pt, bottom=3pt,
        before skip=4pt, after skip=6pt,
        fontupper=\small
    ]
    {\faCodeBranch}\quad{\bfseries\ifIndonesian Kode lengkap\else Full code\fi:}\
    \href{https://github.com/\repoowner/\branchreftemplate}{\texttt{\repoowner/\branchreftemplate}}
    \ifIndonesian
        (repo template, klik \textbf{Use this template} untuk mulai dari kondisi bab ini).
    \else
        (template repo, click \textbf{Use this template} to start from this chapter's state).
    \fi
    \end{tcolorbox}%
}

%% \indexterm{X} writes an index entry that:
%%   - sorts under the lowercased form of X (so chapter uses of \term{widget}
%%     merge with glossary uses of \gitem{Widget} under a single entry),
%%   - displays X with its first character forced to uppercase.
\makeatletter
\newcommand{\indexterm}[1]{%
    \begingroup
        \lowercase{\edef\@@idxkey{#1}}%
        \@@idxcapfirst#1\@nil
        \protected@edef\@@idxentry
            {\noexpand\index{\@@idxkey @\@@idxdisp}}%
        \@@idxentry
    \endgroup
}
\def\@@idxcapfirst#1#2\@nil{%
    \uppercase{\def\@@idxfirst{#1}}%
    \edef\@@idxdisp{\@@idxfirst#2}%
}
\makeatother
\newcommand{\term}[1]{\textit{#1}\protect\indexterm{#1}}
\newcommand{\gitem}[2]{\item[\textit{#1}]\indexterm{#1} #2}

\newcommand{\chaptertag}[1]{%
    \tikz[baseline=(tag.base)]\node[
        fill=black,
        fill opacity=0.08,
        text opacity=1,
        rounded corners=3pt,
        inner xsep=6pt,
        inner ysep=3pt,
        text=black!75,
        font=\sffamily\scriptsize\bfseries
    ] (tag) {#1};%
}

\newcommand{\chapterhero}[5]{%
    \begin{tcolorbox}[
        enhanced,
        boxrule=0pt,
        arc=2pt,
        outer arc=2pt,
        left=14pt,
        right=12pt,
        top=10pt,
        bottom=10pt,
        before skip=2pt,
        after skip=10pt,
        colframe=#1!15!white,
        colback=#1!8!white,
        borderline west={3pt}{0pt}{#1!55!black},
    ]
    {\color{#1!55!black}\sffamily\bfseries\small \ifIndonesian Fokus Bab\else Chapter Focus\fi}\par
    \vspace{3pt}
    {\color{black!85}\fontsize{11}{14}\selectfont #2}\par
    \vspace{8pt}
    \chaptertag{#3}\hspace{6pt}\chaptertag{#4}\hspace{6pt}\chaptertag{#5}
    \end{tcolorbox}%
}

%% ------------------------------------------------------------------------
%% Indeks - imakeidx with two columns, dotted leaders, alphabetic group headers
%% ------------------------------------------------------------------------
\newcommand{\indexgroup}[1]{%
    \par\addvspace{12pt}\nopagebreak
    {\noindent\sffamily\bfseries\large\color{black!70} #1}\par\nopagebreak
    \vspace{1pt}\hrule height 0.4pt\relax
    \par\nopagebreak\vspace{4pt}%
}
\usepackage{imakeidx}
\ifIndonesian
    \makeindex[title=Indeks, intoc, columns=2, columnsep=18pt,
        options={-s indexstyle.ist}]
\else
    \makeindex[title=Index, intoc, columns=2, columnsep=18pt,
        options={-s indexstyle.ist}]
\fi

%% ------------------------------------------------------------------------
%% Teorema dan Definisi - language-aware
%% ------------------------------------------------------------------------
\theoremstyle{definition}
\ifIndonesian
    \newtheorem{definition}{Definisi}[chapter]
    \newtheorem{example}{Contoh}[chapter]
\else
    \newtheorem{definition}{Definition}[chapter]
    \newtheorem{example}{Example}[chapter]
\fi

\theoremstyle{plain}
\ifIndonesian
    \newtheorem{theorem}{Teorema}[chapter]
\else
    \newtheorem{theorem}{Theorem}[chapter]
\fi

%% Konfigurasi nama untuk cleveref (teorema) - language-aware
\ifIndonesian
    \crefname{definition}{definisi}{definisi}
    \Crefname{definition}{Definisi}{Definisi}
    \crefname{example}{contoh}{contoh}
    \Crefname{example}{Contoh}{Contoh}
    \crefname{theorem}{teorema}{teorema}
    \Crefname{theorem}{Teorema}{Teorema}
\else
    % English defaults from cleveref are fine
\fi

%% ------------------------------------------------------------------------
%% Spacing dan Format
%% ------------------------------------------------------------------------
\usepackage{setspace}
\setstretch{1.15}

\setlength{\parindent}{1.2em}
\setlength{\parskip}{0pt}

%% ------------------------------------------------------------------------
%% Chapter dan Section Formatting
%% ------------------------------------------------------------------------
\setkomafont{part}{\fontsize{14}{16}\selectfont\bfseries}
\setkomafont{chapter}{\fontsize{14}{16}\selectfont\bfseries}
\setkomafont{section}{\fontsize{12}{14}\selectfont\bfseries}
\setkomafont{subsection}{\fontsize{11}{13}\selectfont\bfseries}
\setkomafont{subsubsection}{\fontsize{11}{13}\selectfont}
\addtokomafont{paragraph}{\fontsize{11}{13}\selectfont}
\addtokomafont{subparagraph}{\fontsize{11}{13}\selectfont}

%% Compact heading spacing
\RedeclareSectionCommand[beforeskip=0.6cm, afterskip=0.25cm]{chapter}
\RedeclareSectionCommand[
    beforeskip=-1.4ex plus -.8ex minus -.2ex,
    afterskip=0.35ex plus .15ex
]{section}
\RedeclareSectionCommand[
    beforeskip=-1.1ex plus -.8ex minus -.2ex,
    afterskip=0.2ex plus .1ex
]{subsection}
\RedeclareSectionCommand[
    beforeskip=-0.8ex plus -.6ex minus -.2ex,
    afterskip=0.1ex plus .1ex
]{subsubsection}

%% ------------------------------------------------------------------------
%% Miscellaneous
%% ------------------------------------------------------------------------
\usepackage{enumitem}            % Enhanced list environments
\setlist{topsep=1pt, parsep=0pt, itemsep=0.5pt, partopsep=0pt}
\usepackage{csquotes}            % Context sensitive quotation

%% Suppress overfull hbox warnings di batas tertentu
\hfuzz=5pt

%% Let TeX stretch interword space as a last resort instead of producing an
%% overfull box (text spilling past the margin) on tight lines, especially
%% around long inline \texttt-based tokens (\cmd, \phpclass, \route, \file).
\tolerance=1500
\hbadness=1500
\emergencystretch=3em
